\chapter{Background and Related Work}\label{ch:background}

This chapter places the two ideas in the literature and reviews how efficient-training claims have been evaluated. It is organised by topic, and each section ends with what it means for this thesis. References were checked against the original sources as far as stated in \cref{app:reproducibility}; the review is focused rather than exhaustive.

\section{Backpropagation and the credit-assignment problem}
Training a network with parameters $\theta$ means reducing a loss $\mathcal{L}(\theta)$ by gradient steps. Backpropagation \citep{rumelhart1986learning} computes all partial derivatives with one forward and one backward sweep. For a fully connected layer with $n_{\mathrm{in}}$ inputs and $n_{\mathrm{out}}$ outputs and a batch of $B$ examples, the forward pass costs $2Bn_{\mathrm{in}}n_{\mathrm{out}}$ FLOPs, the weight gradient the same amount, and the gradient with respect to the layer input the same amount again (except in the first layer, where the input gradient is not needed). A backward pass therefore costs roughly twice a forward pass, which is why a method that skips part of the backward pass can save a large share of the total (\cref{sec:cost-model}).

Backpropagation has long been criticised as biologically implausible, in particular because it uses the transposed forward weights to send errors backwards (the weight-transport problem). Feedback alignment replaces these weights by fixed random ones \citep{lillicrap2016random}, and direct feedback alignment projects the output error directly to every layer \citep{nokland2016direct}. Decoupled neural interfaces learn the backward signal with a small network \citep{jaderberg2017decoupled}. These methods are mostly studied as \emph{replacements} for backpropagation, and mostly judged by accuracy, not by cost.

\section{Target propagation}
Target propagation computes, for every hidden layer, a target value instead of a gradient, using learned inverse mappings that send a target from the layer above to the layer below \citep{bengio2014auto}. Difference target propagation (DTP) adds a correction term so that the target equals the current activation when the output is already correct \citep{lee2015dtp}. In its simplest form the target for a hidden state $h$ is
\begin{equation}
b \;=\; h + G(\text{target output}) - G(\text{current output}),
\end{equation}
where $G$ maps the output space to the space of $h$ and is trained to invert the forward mapping. \citet{bartunov2018assessing} evaluated DTP, a simplified variant (SDTP), feedback alignment and related methods on MNIST, CIFAR-10 and ImageNet, and reported that these biologically motivated algorithms perform much worse than backpropagation on the larger problems. \citet{meulemans2020theoretical} gave a theoretical framework that relates target propagation to Gauss--Newton-like updates, and \citet{shibuya2023fixed} proposed fixed feedback weights for DTP.

\paragraph{Relation to this thesis.} The loss studied in \cref{ch:study1} is the SDTP-style target applied at a middle layer and used as an \emph{auxiliary} loss together with backpropagation. Prior work uses target propagation to replace backpropagation and judges accuracy; the question here is different (does the extra signal reduce the FLOPs needed to reach a target accuracy?). The thesis does not claim a new loss.

\section{Auxiliary losses and deep supervision}
Adding losses at intermediate layers is an old way to speed up or regularise training: deeply supervised nets attach a classifier to hidden layers \citep{lee2015deeply}, and ladder networks add denoising reconstruction costs at every layer \citep{rasmus2015semi}. Because the meeting loss is also an auxiliary loss, these are the natural controls. The thesis therefore includes a deep-supervision arm (B) and an autoencoder auxiliary (C), and requires that the meeting loss beats a \emph{tuned} baseline, not just an untuned one.

\section{Selecting examples and skipping backward work}\label{sec:related-selection}
A second family of methods spends computation unevenly across examples. Online hard example mining (OHEM) keeps the highest-loss region proposals of a batch for the backward pass in object detection \citep{shrivastava2016ohem}. \citet{loshchilov2016online} rank examples by their most recent loss and sample the next batch with a probability that decays with rank. Importance sampling with a gradient-norm bound reduces the variance of the gradient estimate \citep{katharopoulos2018not}. Ordered SGD builds each update from the top-$q$ losses of a mini-batch and supplies a convergence analysis for it \citep{kawaguchi2020ordered}. RHO-loss prioritises points that are learnable, worth learning and not yet learnt \citep{mindermann2022prioritized}.

Selective-Backprop (SB) \citep{jiang2019selective} is the closest relative of the selective pass studied here. It forwards a batch, computes each example's loss, selects examples with a probability that depends on the percentile of their loss in a history of recent losses (with a hyper-parameter $\beta$ that controls how aggressive the selection is), and back-propagates only the selected ones. The reported motivation is that the backward pass is the more expensive part and that many examples contribute little to learning, and the reported result is a reduction in wall-clock time to a target error on deep image models. SB's published gains are in wall-clock time on deep convolutional networks; this thesis counts FLOPs in a small MLP and uses a reimplementation of the selection rule (\cref{sec:sb}). A direct comparison with the published numbers is therefore not possible, and none is made.

\paragraph{Relation to this thesis.} The selective pass of \cref{ch:methods} is a deterministic top-$k$ variant of this family with a very small number of design choices (kept fraction and warm-up). It is \emph{not} new as an idea. What the thesis adds is a compute-matched, pre-registered test of it against a tuned baseline, a batch-size control, a random-rows control, and a comparison with a tuned SB reimplementation in which every extra forward pass is charged.

\section{Evaluating efficient training}\label{sec:related-eval}
Several studies show how fragile comparisons of training methods are. \citet{schmidt2021crowded} benchmarked fifteen optimisers across many problems and found that tuning protocol matters at least as much as the choice of optimiser. \citet{kaddour2023notrain} compared efficient-training algorithms for transformer language models (including layer stacking, layer dropping and Selective-Backprop) at a fixed budget, evaluated by a ``reference system time'' that makes runs on different hardware comparable, and found that the gains over a well-tuned baseline with a fully decayed learning-rate schedule vanish at the budgets they studied. \citet{dehghani2022efficiency} show that cost indicators (parameter count, FLOPs, speed) can contradict each other, so a single indicator can mislead. \citet{dahl2023benchmarking} propose a benchmark that measures time to a fixed validation target, which is the same style of target-based metric used in this thesis.

Reliability also depends on statistics. \citet{bouthillier2021accounting} show that benchmark conclusions depend on sources of variance (initialisation, data order, hyper-parameter optimisation) that single-seed comparisons ignore. \citet{dodge2019show} argue for reporting the compute budget used for tuning. \citet{bergstra2012random} show that random search is a strong baseline for hyper-parameter optimisation, which is the search used here. Pre-registration, in which hypotheses and analysis plans are fixed before the data are collected, is well established in other empirical sciences \citep{nosek2018preregistration}; this thesis uses a lightweight form of it based on git tags (\cref{sec:prereg}).

\paragraph{Relation to this thesis.} The protocol of \cref{ch:protocol} follows these lessons at small scale: a target-based cost measure, the same tuning budget for every arm, paired seeds with bootstrap intervals, an explicit check of how much re-tuning the baseline changes the result, and locked hypotheses. The finding of \citet{kaddour2023notrain} that a strong baseline removes the apparent gain is a warning that the positive results of \cref{ch:study2,ch:study3} must be read with the limits of \cref{sec:limitations} in mind.

\section{Summary}
Target propagation has been studied as a replacement for backpropagation, not as a cheap auxiliary signal. Selective backward passes have been studied in wall-clock time on deep models, not in counted FLOPs with a batch-size control. Efficient-training claims have been shown to depend on baseline tuning and cost metric. The thesis sits at the intersection: two small, fully controlled tests, with the protocol as the main product.
